% year: תשפ"ב
% type: מבחן
% semester: א
% moed: א
% number: 2א
% points: 8
% categories: func-analysis, func-limits
% source: exams/מבחנים/תשפב/מועד א תשפב.pdf

\begin{question}
(8 נק') מצאו את כל האסימפטוטות של הפונקציה
\[ f(x)=\frac{2x^2-3x+1}{x} \]
\end{question}

\begin{hint}
חלקו את המונה ב-$x$: $f(x)=2x-3+\frac1x$.
\end{hint}

\begin{solution}
תחום ההגדרה: $x\neq0$. לכל $x\neq 0$:
\[ f(x)=2x-3+\frac{1}{x}. \]
\textbf{אסימפטוטה אנכית:} המונה ב-$x=0$ שווה $1\neq0$, ולכן
\[ \lim_{x\to0^+}f(x)=+\infty,\qquad \lim_{x\to0^-}f(x)=-\infty, \]
כלומר $x=0$ אסימפטוטה אנכית. זו הנקודה היחידה מחוץ לתחום, ו-$f$ רציפה בכל נקודה אחרת, לכן אין אסימפטוטות אנכיות נוספות.

\textbf{אסימפטוטה משופעת/אופקית:} נחשב
\[ m=\lim_{x\to\pm\infty}\frac{f(x)}{x}=\lim_{x\to\pm\infty}\left(2-\frac3x+\frac1{x^2}\right)=2, \]
\[ n=\lim_{x\to\pm\infty}\big(f(x)-2x\big)=\lim_{x\to\pm\infty}\left(-3+\frac1x\right)=-3. \]
לכן $y=2x-3$ אסימפטוטה משופעת גם ב-$+\infty$ וגם ב-$-\infty$. כיוון ש-$f(x)\to\pm\infty$ כאשר $x\to\pm\infty$, אין אסימפטוטה אופקית.

\textbf{תשובה:} $x=0$ (אנכית) ו-$y=2x-3$ (משופעת בשני הכיוונים).
\end{solution}
